\section*{Capítulo \ref{ch:b3:microbiomes} --- Microbiomas y simbiosis}

\begin{solution}{exo:b3:microbiomes:1}
\omterm{def:b3:microbiomes:symbiosis}{Mutualismo}: ganan los dos --- leguminosa y rizobio, coral y alga.
\omterm{def:b3:microbiomes:symbiosis}{Comensalismo}: gana uno y el otro no se ve afectado --- la mayoría de las
bacterias intestinales que se alimentan de lo que el hospedador no puede
digerir. Parasitismo: uno gana a costa del otro --- \emph{Wolbachia}
matando embriones macho. Una misma pareja puede cambiar: una comensal del
intestino que cruza la pared se convierte en patógena, un rizobio que
deja de fijar se convierte en parásito del nódulo, y un alga bajo estrés
térmico daña al coral que la alojaba.
\end{solution}

\begin{solution}{exo:b3:microbiomes:2}
$H = -(0.5\ln 0.5 + 0.3\ln 0.3 + 0.1\ln 0.1 + 2\times 0.05\ln 0.05) =
0.347 + 0.361 + 0.230 + 0.300 = 1.24$; $e^{H} = 3.4$ especies efectivas.
$D = 0.25 + 0.09 + 0.01 + 0.0025 + 0.0025 = 0.355$; $1/D =
2.8$.
\end{solution}

\begin{solution}{exo:b3:microbiomes:3}
El estudio del 16S informa de quién está ahí, hasta el género
aproximadamente, y de sus recuentos relativos de lecturas; no dice nada
de los genes que llevan, se pierde hongos y \omterm{def:b3:virology:virus}{virus} (que exigen otros
\omterm{def:b3:genetic-engineering:pcr}{cebadores}) y lo sesgan el número de copias y los \omterm{def:b3:genetic-engineering:pcr}{cebadores}. La
metagenómica de escopetazo informa de genes, de funciones y, con
suficiente profundidad, de genomas y cepas enteros; cuesta más por
muestra, la anega el ADN del hospedador en las muestras de tejido y sigue
sin poder decir qué genes se expresan ni qué células están vivas.
\end{solution}

\begin{solution}{exo:b3:microbiomes:4}
Fermentar la fibra a \omterm{def:b3:microbiomes:gut}{ácidos grasos de cadena corta} que alimentan al colon
y al hospedador; sintetizar vitamina~K y vitaminas del grupo B;
transformar ácidos biliares y fármacos; la \omterm{def:b3:microbiomes:gut}{resistencia a la colonización}
frente a patógenos; y educar al sistema inmunitario. Alteración: la
colitis por \emph{Clostridioides difficile} tras los \omterm{def:b3:bacteriology:antibiotics}{antibióticos}
(también implicadas: la enfermedad inflamatoria intestinal y la
obesidad).
\end{solution}

\begin{solution}{exo:b3:microbiomes:5}
$\binom{20}{5} = \num{15504}$. Especie con $10$: $1 - \binom{10}{5}/
\num{15504} = 1 - 252/\num{15504} = 0.984$; con $6$: $1 - 2002/\num{15504}
= 0.871$; con $3$: $1 - 6188/\num{15504} = 0.601$; con $1$: $1 -
\num{11628}/\num{15504} = 0.25$. $E[S_{5}] = 2.7$ especies. Cobertura: una
singletona, $1 - 1/20 = 0.95$.
\end{solution}

\begin{solution}{exo:b3:microbiomes:6}
$400\times 10^{11} = 4\times 10^{13}$ bacterias, unos \qty{40}{g}. Genes:
$1000$ especies $\times$ \num{4000}, menos el solapamiento --- del orden
de $3\times 10^{6}$ genes distintos, más de cien veces los \num{20000}
humanos.
\end{solution}

\begin{solution}{exo:b3:microbiomes:7}
El oxígeno destruye la \omterm{def:b3:microbiomes:nodule}{nitrogenasa}, pero los bacteroides han de respirar
para producir dieciséis ATP por N$_{2}$. La corteza del nódulo forma una
barrera de difusión que limita la entrada de oxígeno, y la
\omterm{def:b3:microbiomes:nodule}{leghemoglobina}, a concentración milimolar, une el oxígeno que entra de
modo que la concentración libre se mantiene cerca de \qty{10}{nM}
mientras el flujo hacia los bacteroides es alto. El rosado es
oxileghemoglobina; un nódulo verde ha degradado su \omterm{def:b3:microbiomes:nodule}{leghemoglobina} ---
está senescente y ya no fija.
\end{solution}

\begin{solution}{exo:b3:microbiomes:8}
Fallan los dos: ni nódulos ni \omterm{def:b3:microbiomes:mycorrhiza}{micorriza arbuscular}, porque el pico de
calcio es la señal compartida de la \omterm{def:b3:microbiomes:mycorrhiza}{vía común de la simbiosis}. La vía es
por tanto anterior a la nodulación --- se construyó para la \omterm{def:b3:microbiomes:symbiosis}{simbiosis}
con hongos hace cuatrocientos millones de años y las leguminosas la
reclutaron después para los \omterm{def:b3:microbiomes:nodule}{rizobios}.
\end{solution}

\begin{solution}{exo:b3:microbiomes:9}
Kiers expuso algunos nódulos de soja a una atmósfera de argón y oxígeno
en la que los \omterm{def:b3:microbiomes:nodule}{rizobios} no podían fijar nitrógeno, mientras el resto de
los nódulos de la planta fijaban con normalidad; la planta cortó el
suministro de oxígeno a los nódulos que no fijaban y sus \omterm{def:b3:microbiomes:nodule}{rizobios} se
reprodujeron aproximadamente la mitad. Sin esa discriminación, los
\omterm{def:b3:microbiomes:nodule}{rizobios} que se saltaran la costosa fijación se reproducirían más deprisa
que los fijadores, se propagarían por la población y las plantas
acabarían alojando bacterias que no dan nada: el \omterm{def:b3:microbiomes:symbiosis}{mutualismo} se
derrumbaría.
\end{solution}

\begin{solution}{exo:b3:microbiomes:10}
$\binom{N-N_{i}}{n}/\binom{N}{n} = \prod_{j=0}^{n-1}(N - N_{i} - j)/(N -
j)$, con cada factor menor que $1$; pasar de $n$ a $n + 1$ multiplica por
un factor más menor que $1$, de modo que la probabilidad de no verla baja
y la de presencia sube con $n$; en $n = N$ el numerador $\binom{N -
N_{i}}{N}$ es cero y todas las especies están presentes, lo que da $S$.
Para $N_{i},
n \ll N$ cada factor vale aproximadamente $1 - N_{i}/N$, y el producto,
$(1 -
N_{i}/N)^{n} \approx e^{-nN_{i}/N} \approx (1 - n/N)^{N_{i}}$ --- la
aproximación de muestreo con reposición.
\end{solution}

\begin{solution}{exo:b3:microbiomes:11}
Un macho es un callejón sin salida para un simbionte que solo viaja en
los huevos, de modo que todo rasgo que convierta machos en hembras, los
mate en favor de sus hermanas o perjudique a las hembras no infectadas
eleva la transmisión del simbionte. Incompatibilidad citoplasmática: el
esperma de los machos infectados se modifica de manera que el cigoto
muere salvo que el huevo lleve la misma \emph{Wolbachia}, que lo rescata.
Las hembras no infectadas pierden por tanto la descendencia que conciben
con machos infectados, mientras que las infectadas no pierden ninguna; la
ventaja crece con la frecuencia de infección, de modo que por encima de
un umbral la infección se propaga hasta la fijación aunque las hembras
infectadas pongan algunos huevos menos --- el coste queda superado por la
fracción de nidadas de hembras no infectadas que se destruye.
\end{solution}

\begin{solution}{exo:b3:microbiomes:12}
El endosimbionte pasa por un cuello de botella de unas pocas células en
cada generación del hospedador y nunca recombina con extraños: la deriva
fija mutaciones levemente deletéreas, deleciones incluidas, y no hay
fuente de genes de repuesto (el trinquete de Muller). Los genes cuyos
productos aporta el hospedador, o que sirven para una vida libre ---
envoltura, motilidad, regulación, casi toda la reparación ---, no se
echan de menos y se van primero; la maquinaria de traducción y los genes
que fabrican lo que el hospedador necesita (los aminoácidos esenciales en
\emph{Buchnera}) se conservan más tiempo. Los parientes de vida libre
tienen poblaciones enormes, recombinación y entrada de genes, y la
selección mantiene sus genomas. Invertirlo exigiría genes nuevos de fuera
--- imposible en aislamiento ---, de modo que los hospedadores sustituyen
en su lugar al simbionte erosionado por otro nuevo, como han hecho varios
linajes de insectos.
\end{solution}

\begin{solution}{pb:b3:microbiomes:1}
\textbf{1.} $400\times 10^{11} = 4\times 10^{13}$ bacterias, \qty{40}{g}.
\textbf{2.} Un genoma de \num{4000} genes mide unas \qty{4}{Mb},
\qty{4.4}{fg}: $4\times 10^{13}\times 4.4\times 10^{-15} = \qty{0.18}{g}$
de ADN bacteriano, frente a
$3\times 10^{13}\times 6.4\times 10^{-12} = \qty{0.19}{g}$ de ADN humano
--- aproximadamente igual.
\textbf{3.} $1000\times 0.7\times 4000 + 0.3\times 4000 \approx 2.8\times
10^{6}$ genes distintos, unas $140$ veces los del genoma humano.
\textbf{4.} $30\times 10 = \qty{300}{mmol}$; $\times 0.9 = \qty{270}{kJ}$;
el \qty{3}{\%} de la dieta.
\textbf{5.} No --- un \qty{3}{\%} no explica un \qty{30}{\%}. Los ratones
libres de gérmenes absorben además con menos eficiencia, tienen alteradas
las hormonas intestinales, el almacenamiento de grasa y la termogénesis,
y otro apetito: el \omterm{def:b3:microbiomes:symbiosis}{microbioma} actúa sobre la fisiología del hospedador y
no solo como proveedor de combustible.
\textbf{6.} División compensada por la evacuación: una población estable,
con cada célula reemplazada a diario. Tras una matanza del \qty{99.9}{\%}
quedan $4\times 10^{10}$ y diez duplicaciones restauran los números en
unos diez días --- pero dominan las supervivientes y las que crecen más
deprisa, la composición cambia (una \omterm{def:b3:microbiomes:gut}{disbiosis}) y en ese hueco puede
asentarse un invasor como \emph{C. difficile}.
\textbf{7.} Las especies minoritarias se reparten $0.4$ entre $177$:
$p = 0.00226$ cada una. $H
= -(0.3\ln 0.3 + 0.2\ln 0.2 + 0.1\ln 0.1 + 0.4\ln 0.00226) = 0.361 +
0.322 + 0.230 + 2.437 = 3.35$; $e^{H} \approx 28$.
\textbf{8.} $D = 0.09 + 0.04 + 0.01 + 177\times (0.00226)^{2} = 0.141$;
$1/D = 7.1$. El \omterm{def:b3:microbiomes:diversity}{índice de Simpson} eleva al cuadrado las abundancias, de
modo que lo deciden las tres especies dominantes y las raras apenas
cuentan; el de Shannon da más peso a las especies raras.
\textbf{9.} $1 - 60/2000 = 0.97$: unas $30$ lecturas de cada mil
pertenecen a especies aún no vistas.
\textbf{10.} La riqueza sube con el tamaño de la muestra mientras sigan
apareciendo especies raras; \num{10000} lecturas alcanzan especies que
\num{2000} pierden. Rarefáciase la muestra grande a \num{2000} lecturas
con el teorema (o por submuestreo) --- debería dar unas $180$ si es la
misma comunidad --- y compárense a igual profundidad.
\textbf{11.} $N_{i} = 1$: $1 - \binom{1999}{1000}/\binom{2000}{1000} =
1 - (2000 - 1000)/2000 = 0.5$. $N_{i} = 5$: unos $1 - 0.5^{5} = 0.97$.
\textbf{12.} $e^{2.0} = 7.4$ especies efectivas; $D \ge 0.5^{2} = 0.25$,
alrededor de $0.27$: una comunidad con una especie dominante y unas pocas
decenas de especies menores --- el intestino derrumbado y de baja
diversidad en el que se afianza \emph{C.~difficile}.
\textbf{13.} $150\times 8 = \qty{1200}{kg}$ de azúcar por hectárea, el
\qty{10}{\%} de las \qty{12}{t} fotosintetizadas.
\textbf{14.} $150\,000/28 = 5360$ mol de N$_{2}$; $16\times 5360 =
\num{85700}$ mol de ATP.
\textbf{15.} $\num{85700}/30 = 2860$ mol de glucosa, \qty{514}{kg} ---
cerca del \qty{43}{\%} del azúcar pagado. El resto construye y mantiene
los nódulos y los bacteroides, aporta como reductor los ocho electrones
por N$_{2}$ y asimila el amoníaco.
\textbf{16.} $1200\times 16 = \qty{19200}{MJ}$ para \qty{150}{kg} de N:
\qty{128}{MJ/kg}, casi cuatro veces los \qty{35}{MJ/kg} de la fábrica ---
pero pagados con luz solar, en el campo y sin combustible.
\textbf{17.} Fijadores, $0.9\times 1 = 0.9$; tramposos, $0.1\times 0.5 =
0.05$; los tramposos son el $0.05/0.95 = 5.3\,\%$ --- a la mitad en una
sola temporada.
\textbf{18.} Tramposos, $0.1\times 1.2^{10} = 0.62$, frente a fijadores
$0.9$: el \qty{41}{\%} tras diez temporadas, camino del dominio. Las
sanciones vuelven la selección contra el tramposo; sin ellas el
\omterm{def:b3:microbiomes:symbiosis}{mutualismo} se erosiona.
\textbf{19.} $10^{6}\times \qty{20}{pg} = \qty{20}{\micro g}$ de carbono
por cm$^{2}$ y día; \qty{18}{\micro g} para el coral.
\textbf{20.} Respiración, \qty{15}{\micro g}: un excedente de
\qty{3}{\micro g} por cm$^{2}$ y día para el crecimiento, la matriz
orgánica de la calcificación, el moco y los gametos.
\textbf{21.} Con el \qty{5}{\%} de las algas, \qty{0.9}{\micro g} al día
frente a \qty{15}{\micro g} respirados: un déficit de \qty{14}{\micro g}
al día; \qty{300}{\micro g} de reservas duran unas tres semanas.
\textbf{22.} Cuatro semanas a $+2$ es peor: el daño a los fotosistemas de
las algas sube abruptamente, no linealmente, con la temperatura. La
medida funciona de todos modos como aviso porque el blanqueamiento
integra el estrés acumulado y los dos perfiles difieren en menos que la
incertidumbre del umbral; es un índice empírico, calibrado con
blanqueamientos observados.
\textbf{23.} Un coral estresado puede desplazarse hacia especies de algas
más tolerantes al calor ya presentes en poca abundancia, o adquirirlas
del agua tras el blanqueamiento; las algas tolerantes dan menos carbono,
de modo que el coral crece más despacio, y la tolerancia ganada es de uno
o dos grados --- menos que el calentamiento proyectado ---, mientras que
el ritmo del calentamiento adelanta a la adaptación de los propios
corales.
\textbf{24.} $10^{10}$ cm$^{2}$ $\times$ \qty{20}{\micro g} $=
\qty{200}{kg}$ de carbono al día, unas \qty{73}{t} al año.
\textbf{25.} Alrededor del \qty{3}{\%} de la dieta procedente de la
fermentación del \omterm{def:b3:microbiomes:symbiosis}{microbioma}; un \qty{97}{\%} de cobertura de la muestra
de secuenciación; y un \qty{10}{\%} de la fotosíntesis pagado por el
nitrógeno del cultivo de judías.
\end{solution}
